% !TeX root = ../main.tex

\section{Método de Van Noort y método desarrollado}
A continuación, se expone el método original desarrollado en este trabajo para aislar y extraer el campo magnético a partir de las señales observables de Stokes $V$ e $I$. El núcleo de esta aportación se fundamenta en la teoría matemática propuesta por Van Noort \cite{vannoortSpatiallyCoupledInversion2012}, la cual ha sido adaptada e implementada específicamente para resolver el problema de la convolución que nos ocupa.

El marco teórico original introducido por Van Noort se articula del siguiente modo. Supóngase un conjunto de datos observacionales $d$ y un modelo analítico predictivo $F(x)$ dependiente de un vector de parámetros $x$ (en el contexto de este estudio, el modelo corresponde a la aproximación de campo débil y el parámetro es el campo magnético longitudinal). Dada una estimación inicial de los parámetros, denotada como $x_0$, el vector residual o error se define como:

\begin{equation}
    \label{eq:beta}
    \beta = d -F(x_0)
\end{equation}

Para optimizar esta estimación inicial, se busca determinar un vector de corrección $\nu$, de manera que la nueva aproximación actualizada se exprese como $x_0 + \nu$.

Dado que, según la formulación original, la función $F(x)$ es típicamente compleja y no lineal, se procede a linealizarla mediante una expansión en serie de Taylor de primer orden:

\begin{equation}
    F(x_o + \nu) \simeq F(x_0) + J \nu
\end{equation}

donde $J$ es la matriz Jacobiana del sistema. 

La convergencia ideal del algoritmo exige que la nueva estimación concuerde analíticamente con los datos observados, anulando el residuo:

\begin{equation}
    \begin{split}
        d - F(x_o + \nu) &= 0\\
        d - (F(x_0) + J \nu) &= 0\\
        J \nu = \beta
    \end{split}
\end{equation}

donde $\beta = d - F(x_0)$. La dificultad analítica subyacente en el tratamiento de datos experimentales radica en que el número de observaciones supera ampliamente al de parámetros libres a ajustar. En consecuencia, $J$ adquiere la estructura de una matriz rectangular, configurando un sistema algebraico sobredeterminado. 

Para abordar la resolución de este sistema, se implementa una formulación de mínimos cuadrados. Siguiendo la derivación del autor original, la optimización cuadrática del residuo conduce al siguiente sistema de ecuaciones normales:

\begin{equation}
    \label{eq:sistemaResolver}
    J^T J \nu = J^T \beta
\end{equation}

Esta reformulación algebraica permite calcular la corrección, dado que el producto matricial $J^T J$ genera una matriz cuadrada e invertible. Además, la arquitectura del método integra restricciones físicas intrínsecas del sistema para favorecer la convergencia hacia soluciones físicamente viables.

\subsection{Adaptación de la metodología al problema de la convolución polarimétrica}
La matriz $J$ representa el Jacobiano del sistema físico. En este TFG, el modelo directo consta de dos pasos: el cálculo local de la señal de polarización bajo la aproximación de campo débil y su posterior degradación espacial al convolucionar con la PSF. Como cada parámetro $x_i$ está asociado a un píxel de la imagen original, el tamaño del problema bidimensional se dispara. Para imágenes de 1400x1400 píxeles, calcular y almacenar explícitamente la matriz Jacobiana resulta inviable por falta de memoria RAM.

En este punto es donde reside la principal optimización de este método. En lugar de construir $J$ explícitamente como una matriz gigantesca, podemos tratarla matemáticamente como un operador lineal. Esta linealidad garantiza la existencia de su operador adjunto (o transpuesto) $J^T$, lo que permite resolver el sistema (\ref{eq:sistemaResolver}) de forma iterativa sin necesidad de almacenar la matriz en memoria.

Para demostrar esta linealidad, partimos de la aproximación de campo débil (\ref{eq:campoDebil}) y agrupamos todos los términos que no dependen del campo magnético en un único factor espacial $K(x, y, \lambda)$:

\begin{equation}
    V = - C \bar{g} \lambda_B^2 \frac{dI_0}{d \lambda} B_{\parallel} = K(x, y, \lambda) B_{\parallel}
\end{equation}

De este modo, la variación de los datos observados respecto al parámetro libre (el campo $B_{\parallel}$) equivale a multiplicar por el factor $K$ y, a continuación, aplicar la convolución espacial. Como tanto la derivada como la convolución son operadores lineales, el Jacobiano se puede expresar funcionalmente como:

\begin{equation}
    \label{eq:demosLinealidad}
    J = \frac{\partial V_{obs}}{\partial B_{\parallel}} = \frac{\partial}{\partial B_{\parallel}} \left[(K \cdot B_{\parallel}) * P \right] = \left[\frac{\partial}{\partial B_{\parallel}} (K \cdot B_{\parallel})\right] * P = K * P
\end{equation}

Al confirmar que $J$ es un operador lineal, queda garantizada la existencia de su operador traspuesto $J^T$, definido mediante el producto escalar estándar $\langle \cdot , \cdot \rangle$. Para cualquier par de funciones $x$ e $y$, el traspuesto debe cumplir la identidad:

\begin{equation}
    \label{eq:defTraspuesta}
    \langle Jx , y \rangle = \langle x , J^T y \rangle
\end{equation}

Observando la ecuación (\ref{eq:demosLinealidad}), vemos que $J$ es la composición de dos operaciones: un operador de convolución $C_p$ y uno de multiplicación diagonal $M_k$. Como la multiplicación por un escalar es autoadjunta, su traspuesta es ella misma. Por tanto:

\begin{equation}
    \label{eq:defTraspuesConjunta}
    J^T = (C_p M_k)^T = M^T_k C_p^T = M_k C_p^T
\end{equation}

Por otro lado, es conocido que el adjunto de una convolución equivale a convolucionar con el núcleo invertido espacialmente (rotado $180^\circ$), que denotaremos como $\tilde{P}$. Con todo esto, la acción del operador Jacobiano traspuesto sobre un mapa de campo magnético queda definida en su forma definitiva:

\begin{equation}
    \label{eq:jacobianaTras}
    J^T B_{\parallel} = K \cdot (\tilde{P} \ast B_{\parallel})
\end{equation}

\subsection{Resolución numérica: Subespacios de Krylov y Gradiente Conjugado}

Sustituyendo el operador directo (\ref{eq:demosLinealidad}) y su traspuesto (\ref{eq:jacobianaTras}) en el sistema de ecuaciones normales (\ref{eq:sistemaResolver}), el problema se reduce a resolver un sistema lineal de la forma $Ax = b$. 

El esquema computacional del método es el siguiente:

\begin{figure}[H]
    \centering
    \begin{tikzpicture}[
        node distance=0.8cm,
        font=\small,
        block/.style={rectangle, draw, fill=purple!10, text width=6.5cm, text centered, rounded corners, minimum height=1cm},
        arrow/.style={-{Stealth[scale=1.0]}, thick}
        ]
        
        \node[block] (step1) {\textbf{1. Inicialización:} \\ Estimación inicial $B_{\parallel}^0$};
        \node[block, below=of step1] (step2) {\textbf{2. Modelo sintético:} \\ Señal $V' = J B_{\parallel}^0$};
        \node[block, below=of step2] (step3) {\textbf{3. Residuo:} \\ $\beta = V_{\mathrm{obs}} - V'$};
        \node[block, below=of step3] (step4) {\textbf{4. Corrección:} \\ Resolver $J^T J \Delta B_{\parallel} = J^T \beta$};
        
        \draw[arrow] (step1) -- (step2);
        \draw[arrow] (step2) -- (step3);
        \draw[arrow] (step3) -- (step4);
    \end{tikzpicture}
    \caption{Diagrama de flujo del esquema computacional del método de inversión.}
    \label{fig:flowchart_metodo}
\end{figure}

Este proceso da lugar a la siguiente ecuación de optimización:

\begin{equation}
    \label{eq:sistemaEcReal}
    J^T J \Delta B_{\parallel} = J^T \beta
\end{equation}

En la notación algebraica estándar, la matriz del sistema es $A = J^T J$, la incógnita a despejar es la corrección $x = \Delta B_{\parallel}$, y el término independiente es $b = J^T \beta$.

La resolución de este sistema proporciona el mapa bidimensional de correcciones $\Delta B_{\parallel}$, que se suma a la estimación inicial $B_{\parallel}^0$ para generar el mapa de campo magnético actualizado: $B_{\parallel}^f = B_{\parallel}^0 + \Delta B_{\parallel}$.

Dado que calcular y almacenar explícitamente la matriz $A$ sigue siendo inviable, plantear el problema mediante operadores lineales permite emplear métodos iterativos basados en subespacios de Krylov. La principal ventaja de esta familia de algoritmos es que solo requieren evaluar el producto de la matriz por un vector (en este caso, aplicar secuencialmente los operadores sobre la imagen), evitando por completo calcular y guardar los coeficientes de la matriz.

Dentro de los métodos de Krylov, y dado que nuestra matriz $J^T J$ es por construcción simétrica y definida positiva, el algoritmo idóneo para la resolución es el método del Gradiente Conjugado (CG). Este método construye iterativamente direcciones de búsqueda ortogonales (conjugadas) en el subespacio, minimizando el error de forma óptima en cada iteración. Al ser un algoritmo ampliamente estandarizado, su integración resulta directa empleando bibliotecas de cálculo científico convencionales.